% Shared module: the outer loop over operator classes.
\subsection{The outer loop: choosing the class of operators}\label{sec:alg-loop}
A \emph{structure} $\Mfit=((m_1,d_1),\dots,(m_N,d_N))$ fixes the number $N$ of factors and, for each,
its degree $m_j$ in $\tau$ and $d_j$ in $\xi$; a factor is monic in $\tau$. For the running example the true structure is
$\Mfit=((1,1),(1,2))$: one factor of degree one in $\tau$ and one in $\xi$ ($\tau+c\xi$), and one of degree one in $\tau$
and two in $\xi$ ($\tau-\nu\xi^2$). The \emph{class} $\mathcal C(m,d)$ contains all structures with $\sum_jm_j\le m$ and
every $d_j\le d$. The starting class $\mathcal C(3,4)$ contains $((1,1),(1,2))$, but also, for instance, $((2,2))$ (a single
second-order factor, such as the wave operator $\tau^2-c^2\xi^2$) and $((1,0),(1,1),(1,1))$ (three first-order factors);
it has $90$ structures, and the largest class allowed, $\mathcal C(4,8)$, has $1353$. For each structure the core
(\cref{sec:alg-pencil,sec:alg-structfit}) returns fitted factors and a fit score.

\paragraph{Comparing structures} The structures are compared by the Bayesian information criterion
\begin{equation}\label{eq:bic}
  \mathrm{BIC}(\Mfit)=2n_{\rm m}\log\max\big(\chi^2/(2n_{\rm m}),\,1\big)+n_{\rm par}(\Mfit)\log(2n_{\rm m}),
\end{equation}
where $n_{\rm m}$ is the number of fitted modes (each complex rate gives two real data, hence $2n_{\rm m}$),
$n_{\rm par}(\Mfit)$ is the number of free coefficients of the structure, and
\[
  \chi^2=\sum_{p=1}^{n_{\rm m}}\min\big(r_{j(p)}(p)^2,\,9\big)
\]
measures how badly the fitted factors explain the observed modes, through the distance $r_j(p)$ of each mode to the
factor it is assigned to, in units of its standard deviation (\eqref{eq:vardist} in \cref{sec:alg-structfit}). If the
structure is right and the rate errors are circular complex Gaussian, each $r_j(p)^2$ is the squared modulus of a
standard complex Gaussian, exponentially distributed with mean one, so that, ignoring the cap, $\chi^2$ follows over
noise realizations approximately one half of a chi-squared distribution with $2n_{\rm m}-n_{\rm par}$ degrees of
freedom, and $\chi^2/(2n_{\rm m})$ is close to $\tfrac12$. The first term of~\eqref{eq:bic} penalizes a poor fit and the
second penalizes complexity; because of the floor $\max(\cdot,1)$, misfit is penalized only once it clearly exceeds what
the noise explains, and the criterion never rewards a fit that is better than the noise level, which would only be
fitting noise. The criterion is turned into a
probability over structures,
\begin{equation}\label{eq:bic-post}
  P(\Mfit\mid\text{data})\propto e^{-\mathrm{BIC}(\Mfit)/2},
\end{equation}
so that every structure receives a score that trades its fit against its number of coefficients.

\paragraph{Growing the class} The class is grown rather than fixed in advance. Starting from $(m,d)=(3,4)$, the
algorithm raises $m$ or $d$ by one when (a) most of the posterior mass~\eqref{eq:bic-post} is on structures labelled
\MIS{} or \IDK{} (\cref{sec:alg-labels}), or (b) an omnibus goodness-of-fit test rejects the current fit,
$\log_{10}p<-3$, where $p$ is the $p$-value of the test under the hypothesis that the fitted structure explains the data.
A larger class almost always fits better, simply because it has more degrees of freedom, so it is kept only if two
further checks pass. The \emph{held-out check} fits the modes of the earlier part of the time window and predicts those
of the later part: a structure that captures the dynamics predicts data it was not fitted to, while one that
only improves the fit to the data it has seen is overfitting. The \emph{mass-shift check} asks that the larger class
move posterior mass towards structures labelled \FC. Growth stops at $(4,8)$.

\paragraph{Regrouping} The loop runs on every field sample $u^{(1)},\dots,u^{(\Nsamp)}$, and different samples may give
slightly different answers: two factors in one sample and three in another, say. Coefficients of such answers cannot be
compared directly. The algorithm therefore takes the consensus structure, the one with the largest posterior mass over
all samples, and refits every sample at that structure (\emph{regrouping}); after this step every sample contributes one
fitted factor of each type, and factors from different samples can be compared and clustered (\cref{sec:bayes-cloud}).

This loop is the only ``regularization'' in the method, and it is a regularization of the
\emph{class}, by Occam's razor, not a penalty on coefficients. A Tikhonov or sparsity penalty would be
the wrong tool: it cannot make the interaction vanish (\cref{prop:discont}), and it can push the error of the split
below the unbiased floor $\noise^2/\gN$ (\cref{sec:limits}) only by shrinking the components towards a prior, which the
split layer of \cref{sec:bayes-split} does explicitly and reports.
